% Propietario conservado: /Users/ruben/Documents/New project/output/REVISION_ARGUMENTAL_APERTURAS_CIERRES_20260915/VII/delivery/MANUSCRIPT_VII_ES_EN_COMPLETE/source_es/manuscrito/29_retorno_y_memoria_traslacional.tex
% Artículo VII, recuperación aditiva, revisión 04, 10-09-2026.
% RESULTADO_RECUPERADO: S15, 11c_gravedad_holonomia_rev6.tex:16--126.
% El cociclo canónico c27=(1,18) se conserva como antecedente del cómputo
% de rutas; este fragmento prueba su composición y representación.
% La memoria traslacional finita y la amplitud cosmológica s0 tienen
% tipos distintos. El cociclo no evalúa por sí solo una densidad local.

\section{Retour de phase et représentation de la mémoire accumulée}
\label{vii:vii:sec:retorno-memoria-gravedad}

La chronologie des routes enrichies distingue le retour de position,
le retour d'orientation et le retour de phase. La publication observable
identifie certains états d'arrivée, tandis que le registre cumulatif
conserve la séquence parcourue. La représentation géométrique doit
transporter les deux données conjointement.

\subsection{Cocycle d'inversion et changement de feuillet}
\label{vii:vii:subsec:cociclo-memoria}

Soit \(F_{\rm obs}\) la publication chronologique du transport TPK.
Ses blocs canoniques ont pour hiérarchie
\[
\begin{array}{c|l}
27&\text{retour positionnel avec inversion d'orientation},\\
54&\text{retour positionnel et orientationnel},\\
108&\text{retour complet de la phase observable}.
\end{array}
\]
Le registre compte les inversions d'orientation \(g(a)\) et les
changements effectifs de feuillet \(s(a)\). Pour une route,
\begin{equation}
 c(\gamma)=\sum_{a\subset\gamma}(g(a),s(a)),\qquad
 c(\gamma_2\gamma_1)=c(\gamma_2)+c(\gamma_1).
 \label{vii:vii:eq:cociclo-ruta}
\end{equation}
Le calcul canonique du bloc de longueur \(27\) publie
\(c_{27}=(1,18)\). La composition conserve ces incréments dans
les quatre blocs qui complètent la période observable.
Sur les routes orientées vers l'avant, le registre est cumulatif ;
son extension au groupoïde des routes utilise \(\mathbb Z^2\) et
\(c(\bar\gamma)=-c(\gamma)\).

\begin{proposition}[Mémoire du tour observable complet]
\label{vii:vii:prop:memoria-vuelta-completa}
En écrivant \(\mathcal K_{27}=F_{\rm obs}^{27}\), le relèvement
au registre, pour \(x\in\mathcal O_{108}\) et \(z\in\mathbb Z^2\), satisfait
\[
 \widetilde{\mathcal K}_{27}(x,z)
   =(\mathcal K_{27}x,z+(1,18)),\qquad
 \widetilde{\mathcal K}_{27}^{\,4}(x,z)
   =(x,z+(4,72)).
\]
Après \(N\) tours complets, l'incrément est \(N(4,72)\).
\end{proposition}

\begin{proof}
La projection observable de \(\mathcal K_{27}\) est d'ordre quatre. Le calcul canonique précédent fixe \(c_{27}=(1,18)\) de manière constante sur cette orbite observable. L'additivité de \eqref{vii:vii:eq:cociclo-ruta} additionne les quatre incréments et donne \(c_{108}=4(1,18)=(4,72)\). La répétition applique la même loi de concaténation et produit \(Nc_{108}\). Le retour de la première coordonnée conserve l'incrément de la seconde.
\end{proof}

\subsection{Représentation affine du cocycle}
\label{vii:vii:subsec:representacion-afin-cociclo}

La réalisation discrète utilise un élément \(R_4\) d'ordre quatre dans
\(\operatorname{Spin}^+(1,3)\), un vecteur temporel unitaire \(u\)
fixé par son action vectorielle et une unité de longueur \(\ell_0>0\).
L'échelle \(\ell_0\) conserve sa provenance dans la section dimensionnelle.
Nous définissons
\begin{equation}
 \rho_{\rm C}^{\rm disc}(\gamma)=
 \bigl(R_4^{c_g(\gamma)},\,\ell_0c_s(\gamma)u\bigr)
 \in\operatorname{Spin}^+(1,3)\ltimes\mathbb R^{1,3}.
 \label{vii:vii:eq:representacion-discreta-memoria}
\end{equation}
La puissance agit sur le vecteur au moyen de la représentation vectorielle
associée du groupe de spin.

\begin{theorem}[Composition et amplitude translationnelle du retour]
\label{vii:vii:thm:traslacion-memoria}
L'application \(\rho_{\rm C}^{\rm disc}\) conserve l'identité,
la concaténation et l'inversion. Dans les retours complets,
\begin{equation}
 \rho_{\rm C}^{\rm disc}(\gamma_{108})=(I,72\ell_0u),\qquad
 \rho_{\rm C}^{\rm disc}(\gamma_{108}^{\,N})=(I,72N\ell_0u).
 \label{vii:vii:eq:traslacion-retorno-108}
\end{equation}
L'inversion change le signe de la translation.
\end{theorem}

\begin{proof}
Le produit semi-direct est
\((R,a)(S,b)=(RS,a+Rb)\). Comme \(R_4u=u\),
\[
\begin{split}
 \rho_{\rm C}^{\rm disc}(\gamma_2)
 \rho_{\rm C}^{\rm disc}(\gamma_1)
 &=
 \bigl(R_4^{c_g(\gamma_2)+c_g(\gamma_1)},
 \ell_0(c_s(\gamma_2)+c_s(\gamma_1))u\bigr)\\
 &=\rho_{\rm C}^{\rm disc}(\gamma_2\gamma_1).
\end{split}
\]
Le cocycle de l'identité est nul et celui de l'inverse est son opposé.
Pour le tour complet, on substitue \(c_{108}=(4,72)\) et
\(R_4^4=I\). La concaténation de \(N\) tours maintient la partie
rotationnelle égale à l'identité et additionne leurs translations.
\end{proof}

L'amplitude \(72\ell_0\) est une longueur associée à une route finie. Une réalisation au moyen d'une cotétrade et d'une connexion la transporte vers son holonomie affine. La contribution locale aux équations gravitationnelles s'obtient ensuite au moyen de la courbure de cette connexion, de la variation de l'action matérielle et de la normalisation de sa densité. La représentation conserve ainsi le contenu cumulatif exact qui doit intervenir dans cette réalisation.
